\section{A one-sided test for a nonabelian cubic field}\label{sec:cubic}
Let $E=\mathbb Q(\sqrt[3]2)$ and $M=E(\zeta_3)$. Our target is RH for
$\zeta_E$, already beyond the public theorem. The polynomial $X^3-2$
is Eisenstein at two and has discriminant $-108$. Its three roots give
$\operatorname{Gal}(M/\mathbb Q)=S_3$, with quadratic subfield
$\mathbb Q(\sqrt{-3})$. For $p>3$ write $g_p$ for its Frobenius conjugacy
class, and $r_p$ for the number of roots modulo $p$. The following table
records all three possibilities. Write $\tau$ for the character of the
two-dimensional standard representation.
\begin{center}
\begin{tabular}{lrrrr}
\toprule
Factorization mod $p$ & Class size & $r_p$ & $\operatorname{sgn}(g_p)$ & $\tau(g_p)$\\
\midrule
Three linear factors & 1 & 3 & $1$ & $2$\\
One linear and one quadratic factor & 3 & 1 & $-1$ & $0$\\
One irreducible cubic & 2 & 0 & $1$ & $-1$\\
\bottomrule
\end{tabular}
\end{center}
For example, $7$ and $31$ are both $7\pmod{24}$. The cubes modulo seven are $0,1,-1$,
whereas $4^3\equiv2\pmod{31}$ and all three cube roots of unity lie in $\mathbb F_{31}$.
The two traces are therefore $-1$ and $2$. In general $\tau(g_p)=r_p-1$. The character of the permutation action
on the roots is $1+\tau$. Its Artin $L$-function is $\zeta_E$, so the
standard factor is $F=\zeta_E/\zeta$. In the right half of the critical
strip $F$ is holomorphic using the repository's ordinary RH, and
$F(1)$ is finite and positive by the simple poles of both zeta functions.
This local argument suffices everywhere we use $F$.

\paragraph{Theorem (three splitting types).}
Use Dirichlet GRH as established in the pinned source. Choose any real weights
$w=(u,v,z)$ on the three rows, in order. Put
\[
 A=(u+3v+2z)/6,\quad B=(u-3v+2z)/6,\quad C=(u-z)/3.
\]
Thus $w=A+B\operatorname{sgn}+C\tau$. Set
\[
 S_w(x)=\sum_{3<p\le x}w(g_p)\log p\,(1/p-1/x),\qquad
 R_w(x)=S_w(x)-A\log x-c_w.
\]
In this sum and below $p$ ranges over primes. The constant $c_w$ is fixed by
$R_w(x)\to0$. The following statements classify the tests obtained from these weights.
\begin{enumerate}
\item If $C=0$, then $R_w(x)=O_w(x^{-1/2}\log^2(2x))$ follows from Dirichlet GRH.
\item If $C\ne0$, RH for $\zeta_E$ is equivalent to the following
\emph{one-sided} estimate. For every $\epsilon>0$ there exists
$K_{w,\epsilon}>0$ such that, for all sufficiently large $x$,
\[
 \operatorname{sgn}(C) R_w(x)\le K_{w,\epsilon}x^{-1/2+\epsilon}.
\]
Under RH for $\zeta_E$ one also gets the two-sided error in part 1 for these weights.
\end{enumerate}
For example, just an upper bound for primes at which $X^3-2$ splits
completely suffices. Just a lower bound for primes at which it is
irreducible also suffices. For $w=(2,0,-1)$ the logarithmic main term
vanishes. This last choice measures exactly the part missed by the two
Dirichlet characters in the table.

\paragraph{Constants and direct estimates.}
Let $\theta_w(t)=\sum_{3<p\le t}w(g_p)\log p$. Integration gives
$S_w(x)=\int_1^x\theta_w(t)t^{-2}\,dt$. The fixed-field prime ideal
theorem for $E$, $\mathbb Q(\sqrt{-3})$ and $\mathbb Q$ gives
$\theta_w(t)=At+O_w(t e^{-b\sqrt{\log t}})$ for some $b>0$ and large $t$.
Indeed degree-one primes of $E$ above $p>3$ are counted by $r_p$, and the
quadratic field similarly gives $1+\operatorname{sgn}(g_p)$.
Terms of higher residue degree and higher ideal powers add
$O(\sqrt t\log^2(2t))$ in each fixed field, by summing over $p^j\le t$,
$j\ge2$. A possible fixed exceptional zero contributes $O(t^{\beta_0})$,
which is absorbed for large $t$. Thus
\[
 c_w=\int_1^\infty(\theta_w(t)-At)t^{-2}\,dt,\qquad
 R_w(x)=-\int_x^\infty(\theta_w(t)-At)t^{-2}\,dt .
\]
Under RH for the respective Dedekind zeta function, the prime ideal
theorem has error $O(\sqrt t\log^2(2t))$. For the quadratic field this
follows from $\zeta_{\mathbb Q(\sqrt{-3})}=\zeta L(s,\chi_{-3})$.
The last integral now proves part 1, and the forward assertion in part 2.
For the prime ideal estimate see
\href{https://arxiv.org/abs/1312.4463}{Greni\'e and Molteni, \emph{Explicit versions of the prime ideal theorem}}.

\paragraph{Why the reverse implication needs only one side.}
By linearity and part 1 it suffices to work with $w=\tau$ and an upper
bound. Initially for $\Re s>1$ set
\[
 D(s)=\sum_{p>3}(r_p-1)(\log p)p^{-s}.
\]
The Euler factors give $D(s)=-F'(s)/F(s)-H(s)$ with $H$ holomorphic
on $\Re s>1/2$. Indeed the $j$th power coefficient of an unramified
prime is $\tau(g_p^j)\log p$ and $|\tau|\le2$. The sum for $j\ge2$
converges absolutely locally uniformly there. The two ramified primes
have locally holomorphic corrections, since their local inverse roots are
zero or roots of unity. The prime ideal theorem above gives
$c_\tau=D(1)$ by Abelian passage to the limit, or by integrating
$D(s)=s\int_1^\infty\theta_\tau(x)x^{-s-1}\,dx$ down to $s=1$.

Use the repository's smoothing and take a Mellin transform. Absolute convergence for
$\Re s>1$ gives, by integrating from $p$ to infinity,
\[
 T(s)=\int_1^\infty R_\tau(x)x^{-s}\,dx
      =\frac{D(s)}{s(s-1)}-\frac{D(1)}{s-1}.
\]
The apparent singularity at one is removable. Each zero $\rho$ of $F$ in
$\Re s>1/2$ supplies a simple pole of $T$ with residue
$m_\rho/[\rho(1-\rho)]$, where $m_\rho$ is the multiplicity. In particular
every real pole in that half-plane has \emph{positive} residue, since real
zeros there lie strictly below one.

Fix $1/2<\beta<1$ and take the assumed bound with $\epsilon=\beta-1/2$.
The function $U(x)=K x^{\beta-1}-R_\tau(x)$ is nonnegative beyond a finite
$X$. Its Mellin integral from $X$ to infinity equals
$K/(s-\beta)-T(s)$ up to an entire function. Suppose its abscissa of
convergence $\sigma_c$ exceeds $\beta$. It is at most one since
$R_\tau(x)\to0$. Landau's lemma for a nonnegative function forces a
singularity at the real point $\sigma_c$. Meromorphic continuation here
permits only a simple pole with \emph{negative} residue, and this would make
the integral negative as $s\downarrow\sigma_c$ on the real axis. This
contradiction shows $\sigma_c\le\beta$. The integral is holomorphic
throughout $\Re s>\beta$, so $T$ has no poles there. Taking all $\beta>1/2$
proves that $F$ is zero-free on the right, and hence so is $\zeta_E$ away
from its pole. The Dedekind functional equation gives the left half.
This proves the equivalence. In this argument the prime-power correction
is holomorphic on the \emph{open} half-plane, which is precisely the
domain needed for the conclusion.

\paragraph{An exact limit on what congruence weights detect.}
Even a very large fixed modulus cannot recover the missing trace.
For any complex function $a$ on the units modulo a fixed $q$,
Chebotarev gives
\[
 \lim_{x\to\infty}\frac1{\pi(x)}
 \sum_{\substack{p\le x\\p\nmid 6q}} |\tau(g_p)-a(p\bmod q)|^2
   =1+\frac{1}{\varphi(q)}\sum_{b\in(\mathbb Z/q)^\times}|a(b)|^2 .
\]
To prove this, form $M\mathbb Q(\zeta_q)$. The intersection of the
two fields is abelian over $\mathbb Q$, so it is either $\mathbb Q$ or the
quadratic subfield of $M$. Conditional on a cyclotomic Frobenius, the
$S_3$ component is uniform on the entire group or on a coset of $A_3$.
The trace has mean zero on the even coset ($2,-1,-1$) and on the odd
coset ($0,0,0$). Its unconditional mean square is $(4+2)/6=1$.
Expanding the square proves the formula, because Chebotarev equidistributes
the compatible pairs with their conjugacy multiplicities.
In particular, linear combinations of any finite set of Dirichlet characters
cannot distinguish this trace from zero in their correlations.
This rules out an exact encoding by fixed residue classes. Weighted
estimates exploiting genuinely new arithmetic structure remain possible.

\paragraph{What this would give for Artin's conjecture.}
For base $2$, Hooley's argument uses the fields
$\mathbb Q(\zeta_n,2^{1/n})$ for squarefree $n$. Our $M$ is the instance
$n=3$. Its zeta function factors as
$\zeta_M=\zeta\,L(s,\chi_{-3})F^2$ by the regular representation of $S_3$.
Thus the criterion identifies one concrete missing instance, including
its precise sign and arithmetic weight. Resolving the full primitive-root
problem by this route would require the family and the analytic control
used in its sieve. For the precise conditional theorem see
\href{https://arxiv.org/abs/2505.05601}{Fan and Pollack, \emph{Counting primes with a given primitive root, uniformly}}.
